% AST2000 Part 4: onboard orientation software (pointing, velocity, position).
% Compile: cd "Project 4/paper" && latexmk -pdf part4.tex   (or: python build_papers.py 4 from the repo root)
% Every number written as \val{key} comes from results.tex, made by pro_4.py.
\documentclass[aps,prd,twocolumn,nofootinbib,floatfix,superscriptaddress]{revtex4-2}
\usepackage{../../common/paper/ast2000}

\begin{document}

\title{Part 4: Where Are We, and Which Way Are We Going? Onboard Orientation Software}
\author{Sakarias Frette}
\affiliation{AST2000 Rerun, University of Oslo}
\date{\today}

\begin{abstract}
A spacecraft must find its own state from its instruments.
We build software that does this right after launch.
\begin{itemize}
\item \textbf{Pointing.} We match a camera picture against 360 reference pictures that we project stereographically from a map of the full sky. The tested pointing angles are accurate to \val{pointing-max-err}$^\circ$, and remain within the API's $2^\circ$ tolerance up to pixel noise of $\sigma=\val{pointing-noise-ok}$.
\item \textbf{Velocity.} We combine the Doppler shifts of two reference stars, which gives \val{v-craft-kms}\,km/s relative to the star.
\item \textbf{Position.} We trilaterate from the distances to all planets and the star, by a linear least-squares solution that needs no initial guess.
\end{itemize}
The software agrees with our own launch prediction to \val{mo-pos-diff}\,km and \val{mo-vel-diff}\,m/s, and \code{ast2000tools} verified it (\val{mo-verified}).
Along the way we found that the API defines the Doppler measurement differently from the textbook.
With our star geometry, the textbook reading would be \val{textbook-err}\% off.
\end{abstract}

\maketitle

% =====================================================================
\section{Introduction}
After launch, the spacecraft knows nothing about its state except what its instruments tell it~\cite{ast2000project}.
It has three instruments:
\begin{itemize}
\item a camera, which gives the pointing angle (\cref{sec:pictures});
\item a spectrograph that measures Doppler shifts, which gives the velocity (\cref{sec:doppler});
\item a radar that measures distances, which gives the position (\cref{sec:tri}).
\end{itemize}
The function \code{orient(mission)} combines all three.
Part 5 calls it again and again during the journey.
Solar system seed: \val{seed}.

% =====================================================================
\section{Pictures of the sky}\label{sec:pictures}
\subsection{Stereographic projection}
The sky map is a HEALPix map~\cite{gorski2005} with one colour per direction $(\theta,\phi)$.
A camera with a $70^\circ$ field of view sees a flat picture: a stereographic projection of a patch of the sphere.
For a picture centred at $\theta_0=90^\circ$, $\phi_0$, a point $(X,Y)$ in the picture plane maps back to
\begin{align}
  \theta &= \frac{\pi}{2}-\arcsin\frac{Y\sin\beta}{\rho}, &
  \phi &= \phi_0+\arctan\frac{X\sin\beta}{\rho\cos\beta},
\end{align}
with $\rho=\sqrt{X^2+Y^2}$ and $\beta=2\arctan(\rho/2)$.
The API looks up one pixel per call.
360 full pictures would take $10^8$ calls, so we ported its HEALPix formula to NumPy.
Our lookup agrees with the API's in all 20\,000 tested directions (\val{hp-mismatch} mismatches).
Our picture for $\phi_0=0$ equals the course's sample picture up to its added noise, which has a standard deviation of \val{sample-res-std} (\cref{fig:proj}).

\subsection{Finding the pointing angle}
We compare a picture with reference pictures for $\phi_0=0,1,\ldots,359^\circ$ by the mean squared colour difference.
We take the best whole degree and refine it with a parabola through its neighbours.
The key step is to compare the pictures after averaging \val{block}$\times$\val{block} pixel blocks.
Raw pixels are dominated by sharp stars, and a picture rotated by only half a degree already differs as much as a random one (\cref{fig:robust}, right).
A library of whole degrees would then have no valley to slide into.
The blurred pictures keep the bright galactic band, which changes smoothly with angle, and so give one wide minimum.

\subsection{Results}
The software finds \val{angle-sample0}$^\circ$ for the sample picture, \val{angle-sample200}$^\circ$ for the second sample, and \val{angle-craft}$^\circ$ for the spacecraft's own picture (\cref{fig:pointing}).
On 20 noisy pictures at random angles the worst error is \val{pointing-max-err}$^\circ$.
\Cref{fig:robust} shows the noise margin.
The pointing stays within $2^\circ$ up to pixel noise of $\sigma=\val{pointing-noise-ok}$, far above the noise of \val{sample-res-std} in real pictures.
It fails only once the noise is comparable to the galactic band's contrast.

\widefig{a_stereographic_projection}{The whole sky (top) with the area a $70^\circ$ picture covers, our projection, the sample picture, and their difference.}{fig:proj}
\widefig{b_pointing}{A panorama of the reference pictures along the ecliptic, and the match score against angle for each test picture.}{fig:pointing}
\widefig{b2_pointing_robustness}{Left: pointing error against added pixel noise. Right: match score near the true angle for different block sizes; blurring is what makes a whole-degree library work.}{fig:robust}

% =====================================================================
\section{Velocity from Doppler shifts}\label{sec:doppler}
\subsection{Method}
A source receding at speed $v_r\ll c$ shifts a spectral line by $\Delta\lambda/\lambda_0=v_r/c$, with $\lambda_0=656.3$\,nm (H$\alpha$).
Two reference stars, at $\phi=\val{star-angle-1}^\circ$ and $\val{star-angle-2}^\circ$, recede from the Sun at \val{sun-v1} and \val{sun-v2}\,m/s.
The difference between the shifts seen from the Sun and from the spacecraft gives the spacecraft's speed $v_i$ towards each star.

Turning $(v_1,v_2)$ into $(v_x,v_y)$ has two readings.
\begin{itemize}
\item The textbook projection, $v_i=\vb v\cdot\hat{\vb u}_i$.
\item The basis expansion $\vb v=v_1\hat{\vb u}_1+v_2\hat{\vb u}_2$, which is how \code{ast2000tools} generates its data, as we found by reading its source.
\end{itemize}
The two agree only for perpendicular stars.
Ours are \val{star-sep}$^\circ$ apart, and the textbook reading would be \val{textbook-err}\% off (\cref{fig:doppler}).

\subsection{Results}
Right after launch the spacecraft moves at $(\val{v-craft-x},\val{v-craft-y})\AU/\yr$, which is \val{v-craft-kms}\,km/s.
Self-tests check three cases:
\begin{itemize}
\item the Sun's own shifts give $\vb v=0$;
\item random velocities survive a round trip through simulated shifts;
\item both readings agree for perpendicular stars.
\end{itemize}

The error budget (\cref{fig:dopplererr}) is simple.
Every $10^{-6}$\,nm of wavelength error is \val{ms-per-micro-nm}\,m/s along one star direction.
With the API's basis expansion, the geometry of the stars does not amplify this.
To stay within the API's tolerance of 1\% of $|\vb v|$ (\val{v-tolerance}\,m/s), each shift must be measured to better than \val{doppler-precision}\,nm.

\widefig{c_doppler_velocity}{Left: the velocity built from the two star directions. Right: error of the textbook reading against the angle between the stars.}{fig:doppler}
\widefig{c2_doppler_errors}{Left: velocity error against the precision of the Doppler measurement. Right: the error ellipse for $10^{-4}$\,nm noise.}{fig:dopplererr}

% =====================================================================
\section{Position by trilateration}\label{sec:tri}
\subsection{Method}
At time $t$ the planets' positions $\vb p_i$ are known exactly from Kepler's equation (Part 2), and the star sits at the origin.
The radar gives the distance $d_i$ to every planet and $d_s$ to the star:
\begin{equation}
  |\vb x-\vb p_i|^2=d_i^2, \qquad |\vb x|^2=d_s^2.
\end{equation}
Subtracting the star's equation from each planet's removes $|\vb x|^2$, the only nonlinear term:
\begin{equation}
  2\vb p_i\cdot\vb x = |\vb p_i|^2+d_s^2-d_i^2. \label{eq:tri}
\end{equation}
This gives eight linear equations in two unknowns, solved by least squares in one step, without iteration or an initial guess.

\subsection{Results}
The position right after launch is $(\val{tri-pos-x},\val{tri-pos-y})\AU$, with all residuals below \val{tri-residual}\,km (\cref{fig:tri}).
The radar's errors are proportional to distance.
Propagating them through \cref{eq:tri} gives the error map in \cref{fig:trimap}: about \val{tri-err-here}\,km per $10^{-6}$ relative noise where the spacecraft is.
Of the noise levels we tested, it first leaves the API's 0.01\,AU tolerance at a relative noise of \val{tri-noise-limit}.

\widefig{d_trilateration}{Every distance is a circle around its planet; all of them cross at the spacecraft. Right: error against noise.}{fig:tri}
\widefig{d2_error_map}{Expected position error across the system for $10^{-6}$ relative noise on every distance, from linear error propagation through the least-squares solution.}{fig:trimap}

% =====================================================================
\section{Manual orientation}
With the three parts combined, the spacecraft finds the following (\cref{fig:mo}):
\begin{itemize}
\item pointing angle \val{mo-angle}$^\circ$;
\item velocity $(\val{v-craft-x},\val{v-craft-y})\AU/\yr$;
\item position $(\val{tri-pos-x},\val{tri-pos-y})\AU$.
\end{itemize}
These come from instruments only.
Our launch simulation (Parts 1 and 3) predicted a state that agrees to \val{mo-pos-diff}\,km and \val{mo-vel-diff}\,m/s.
\code{ast2000tools} verified all three (\val{mo-verified}).
The \val{mo-vel-diff}\,m/s difference is the API's Euler-integrated launch against our adaptive one.
Over a four-month journey it would grow to a large miss, which is why Part 5 re-aims from the measured state.

\widefig{e_orientation_summary}{Manual orientation. The spacecraft's picture, our projection at the angle found, and their difference. Below: onboard fix against launch prediction, with the API's tolerances.}{fig:mo}

% =====================================================================
\section{Discussion and conclusion}
Each instrument has a clear error model:
\begin{itemize}
\item the pointing is limited by the contrast of the galactic band;
\item the velocity by the spectrograph's precision;
\item the position by the radar's relative precision and the planets' geometry.
\end{itemize}
In all three cases our margins are large.
The most important finding was the Doppler convention.
Following the assignment's formula would have failed verification with our stars, and only reading the API's source revealed why.

\bibliography{../../common/paper/references}

\appendix
\section{Mission log}
\begin{missionbox}{Part 4 key numbers (seed \val{seed})}
Pointing after launch & \val{mo-angle}$^\circ$ \\
Velocity & (\val{v-craft-x}, \val{v-craft-y})\,AU/yr \\
Position & (\val{tri-pos-x}, \val{tri-pos-y})\,AU \\
vs.\ launch prediction & \val{mo-pos-diff}\,km, \val{mo-vel-diff}\,m/s \\
Worst pointing error (tests) & \val{pointing-max-err}$^\circ$ \\
Textbook Doppler error & \val{textbook-err}\% \\
API verification & \val{mo-verified} \\
\end{missionbox}

\end{document}
